% !TeX root = ../main.tex
\documentclass[12pt,letterpaper,twoside,openany]{report}

% Input encoding (not needed with XeLaTeX)
\usepackage[utf8]{inputenc}

% Layout: allow fine-grained control over margins, spacing,
% section titles, tables of contents, and other formatting details
\usepackage{array,stfloats,wrapfig}
\usepackage{geometry,setspace,titlesec,tocloft}

% Margins: keep footnotes flush and anchored at bottom of page,
% and enumerate each footnote with a number
\usepackage[hang,flushmargin,bottom]{footmisc}
\usepackage{enumitem}

% Figures and Tables: allow tables and figures that span multiple pages
% or require pages to be oriented in landscape mode, as well as tables
% where some cells can span multiple columns or rows
\usepackage{lscape,pdflscape,longtable,rotating,multirow,multicol}

% Figures and Tables: general features
\usepackage{graphicx,xcolor}
\usepackage{epsfig} % Allow use of EPS images
\usepackage{listings}   % Allow syntax highlighting of code snippets
\usepackage{caption}    % Allow caption formatting and spacing control
\usepackage{float}      % Give additional control for arranging figures
\usepackage{subcaption} % Allow 2+ captioned figures/tables in one float
\usepackage{placeins}   % Restrict float position using FloatBarrier
\usepackage{afterpage}  % Defer non-float blocks (e.g. landscape longtables)
\usepackage{comment}    % Allow special comments for debug/development
\usepackage{colortbl}   % Allow coloring of table cells

% Tables: allow cleaner tables with better spacing and formatting
\usepackage{booktabs,threeparttablex}

% Math: allow standard math symbols, (sub)hypotheses for formal
% proof statements, and pseudo-code writing. Note that coding and
% table-related packages can interact with math packages - sometimes
% in necessary ways (e.g. multirow for multi-row equations), but
% sometimes in undesirable ways (e.g. booktabs and amsmath can interact
% to cause extra vertical space in tables).
\usepackage{amsmath,amssymb,amsfonts,ntheorem,mathrsfs}
\usepackage{algorithm,algorithmic} % replace with algorithm2e, if desired
\newtheorem{hyp}{Hypothesis} 
\newtheorem{subhyp}{Hypothesis}
\renewcommand\thesubhyp{\thehyp\alph{subhyp}}

% Physics: allow proper typesetting of units and numbers
% without needing to use awkward math mode syntax
\usepackage{siunitx}
\sisetup{
  mode = match,
  propagate-math-font = true,
  reset-math-version = false,
  reset-text-family = false,
  reset-text-series = false,
  reset-text-shape = false,
  text-family-to-math = true,
  text-series-to-math = true
}

% General Typography
\usepackage{times}       % Use Times New Roman font
\usepackage{textcomp}    % Allow proper typesetting of text symbols
\usepackage{textcase}    % Text-only case transforms (leave math/symbols intact)
%\usepackage{fontspec}    % Use system fonts (for XeLaTeX and LuaLaTeX)
\usepackage{indentfirst} % Indent first paragraph of each section

% Macro support: allow macros that do not consume space in text mode
\usepackage{xspace}

% PDF links: allow clickable links in the PDF output; hide link borders
\usepackage[hidelinks,unicode=true]{hyperref} % backref=page
\usepackage{bookmark}

% Section header formatting:
% - Center-align and uppercase chapter titles
% - Single-space chapter titles and section headers
% - Retain double-spacing between chapter titles and following text
% - Retain double-spacing between section headers and following text
\titleformat*{\section}{\bfseries\fontsize{12}{14}\selectfont\singlespacing}
\titleformat{name=\chapter,numberless}[block]
{\normalfont\normalsize\filcenter\singlespacing}{}{0pt}{\normalsize}
\titleformat*{\subsection}{\fontsize{12}{14}\selectfont\itshape\singlespacing}

% Bibliography
%\usepackage{cite}
%\usepackage[numbers]{natbib}
%\renewcommand*{\bibfont}{\footnotesize}
%\setlength{\biblabelsep}{\labelsep}
\usepackage[
  backend=biber,
  bibstyle=ieee,
  citestyle=numeric-comp,
  sorting=none,
  sortcites=true,
  dashed=true]{biblatex}
\renewcommand*{\bibopenbracket}{}  % Remove brackets around reference no.
\renewcommand*{\bibclosebracket}{}
\DeclareFieldFormat{labelnumber}{\mkbibbrackets{[#1]}}
\addbibresource{references.bib}
\AtEveryBibitem{
  \clearfield{urlyear}  % Access date is hidden
  \clearfield{urlmonth}
  \clearfield{url}      % non-DOI URL is hidden
}
\setlength{\bibitemsep}{12pt}       % Double-space between references
\appto{\bibsetup}{\setstretch{1.0}} % Single-space within each reference
\usepackage{microtype}              % Improves justification

% Call local resources
\input{common/layout.tex}
\input{common/macros.tex}
